% C19: 原模式、原时长，P/S唯一变更为B.pre起点0ms/10ms。只聚焦CPU预处理与高档域轨迹。
\begin{figure}[htbp]
\centering
\resizebox{0.96\textwidth}{!}{%
\begin{tikzpicture}[font=\small,>=Latex,x=0.36cm,y=1cm]
 \node[anchor=west,font=\bfseries] at (-5.1,4.15) {安排P：并行};
 \node[anchor=west,font=\bfseries] at (-5.1,1.65) {安排S：接力};
 \foreach \y/\lbl in {3.45/{CPU1：A.pre},2.65/{CPU2：B.pre},1.00/{CPU1：A.pre},.20/{CPU2：B.pre}}{
  \node[anchor=east,font=\scriptsize] at (-.35,\y) {\lbl};
 }
 \filldraw[draw=blue!70!black,fill=blue!17,thick] (0,3.19) rectangle (10,3.71);
 \filldraw[draw=orange!70!black,fill=orange!18,thick] (0,2.39) rectangle (15,2.91);
 \filldraw[draw=blue!70!black,fill=blue!17,thick] (0,.74) rectangle (10,1.26);
 \filldraw[draw=orange!70!black,fill=orange!18,thick] (10,-.06) rectangle (25,.46);
 \node[font=\scriptsize] at (5,3.45) {H，10 ms};
 \node[font=\scriptsize] at (7.5,2.65) {H，15 ms};
 \node[font=\scriptsize] at (5,1.00) {H，10 ms};
 \node[font=\scriptsize] at (17.5,.20) {H，15 ms};
 \draw[blue!70!black,very thick] (0,1.93)--(15,1.93);
 \draw[orange!75!black,very thick] (0,-.55)--(25,-.55);
 \node[font=\scriptsize,anchor=east] at (-.35,1.93) {CPU域H};
 \node[font=\scriptsize,anchor=east] at (-.35,-.55) {CPU域H};
 \foreach \x in {0,10,15,25,30}{
  \draw[densely dotted,gray!65] (\x,-1.1)--(\x,3.9);
  \node[below,font=\scriptsize] at (\x,-1.16) {\x};
 }
 \draw[->] (0,-1.10)--(33,-1.10) node[right,font=\scriptsize] {ms};
 \node[font=\scriptsize,blue!70!black] at (18,1.90) {15 ms};
 \node[font=\scriptsize,orange!70!black] at (28,-.56) {25 ms};
 \node[align=center,font=\scriptsize] at (14,-2.0) {其余Stage的设备、OPP和起点完全相同；A+C场景没有变化};
\end{tikzpicture}}
\caption{C19数学见证中A+B场景的预处理时序与共享CPU域H档驻留（其余Stage省略）}
\label{fig:ch3-c19-schedules}
\end{figure}
